% OpenStax Calculus Volume 3 -- Section 5.1 Exercises: Double Integrals over Rectangular Regions
% Exercises reproduced from OpenStax, licensed under CC BY 4.0.
% Source: https://openstax.org/books/calculus-volume-3/pages/5-1-double-integrals-over-rectangular-regions
\documentclass{exercises}
\title{OpenStax Calculus Volume 3}
\subtitle{Section 5.1 Exercises: Double Integrals over Rectangular Regions}
\sourceurl{https://openstax.org/books/calculus-volume-3/pages/5-1-double-integrals-over-rectangular-regions}
\author{}
\date{}

\begin{document}
\maketitle


In the following exercises, use the midpoint rule with $m=4$ and $n=2$ to estimate the volume of the solid bounded by the surface $z=f(x,y)$, the vertical planes $x=1$, $x=2$, $y=1$, and $y=2$, and the horizontal plane $z=0$.

\begin{enumerate}[start=1]
\item $f(x,y)=4x+2y+8xy$
\item $f(x,y)=16x^{2}+\frac{y}{2}$
\end{enumerate}

In the following exercises, estimate $\iint_{R}f\left( x,y \right)\,dA$ by using a Riemann sum with $m=n=2$ and the sample points to be the lower left corners of the subrectangles of the partition.

\begin{enumerate}[resume]
\item $f(x,y)=\sin x-\cos y$, $R=[0,\pi]\times [0,\pi]$
\item $f(x,y)=\cos x+\cos y$, $R=[0,\pi]\times [0,\frac{\pi}{2}]$
\item Use the midpoint rule with $m=n=2$ to estimate $\underset{R}{\iint}f(x,y)\,dA$, where the values of the function $f$ on $R=[8,10]\times [9,11]$ are given in the following table. \begin{center}
\begin{tabular}{lccccc}
\hline
 & \multicolumn{5}{c}{y} \\
x & 9 & 9.5 & 10 & 10.5 & 11 \\
8 & 9.8 & 5 & 6.7 & 5 & 5.6 \\
8.5 & 9.4 & 4.5 & 8 & 5.4 & 3.4 \\
9 & 8.7 & 4.6 & 6 & 5.5 & 3.4 \\
9.5 & 6.7 & 6 & 4.5 & 5.4 & 6.7 \\
10 & 6.8 & 6.4 & 5.5 & 5.7 & 6.8 \\
\hline
\end{tabular}
\end{center}
\item The values of the function $f$ on the rectangle $R=[0,2]\times [7,9]$ are given in the following table. Estimate the double integral $\underset{R}{\iint}f(x,y)\,dA$ by using a Riemann sum with $m=n=2$. Select the sample points to be the upper right corners of the subsquares of $R$. \begin{center}
\begin{tabular}{lccc}
\hline
 & $y_{0}=7$ & $y_{1}=8$ & $y_{2}=9$ \\
$x_{0}=0$ & 10.22 & 10.21 & 9.85 \\
$x_{1}=1$ & 6.73 & 9.75 & 9.63 \\
$x_{2}=2$ & 5.62 & 7.83 & 8.21 \\
\hline
\end{tabular}
\end{center}
\item The depth of a children's 4-ft by 4-ft swimming pool, measured at 1-ft intervals, is given in the following table.
\begin{enumerate}[label=(\alph*)]
  \item Estimate the volume of water in the swimming pool by using a Riemann sum with $m=n=2$. Select the sample points using the midpoint rule on $R=[0,4]\times [0,4]$.
  \item Approximate the average depth of the swimming pool.\\ \begin{center} \begin{tabular}{lccccc} \hline  & \multicolumn{5}{c}{y} \\ x & 0 & 1 & 2 & 3 & 4 \\ 0 & 1 & 1.5 & 2 & 2.5 & 3 \\ 1 & 1 & 1.5 & 2 & 2.5 & 3 \\ 2 & 1 & 1.5 & 1.5 & 2.5 & 3 \\ 3 & 1 & 1 & 1.5 & 2 & 2.5 \\ 4 & 1 & 1 & 1 & 1.5 & 2 \\ \hline \end{tabular} \end{center}
\end{enumerate}
\item The depth of a 3-ft by 3-ft hole in the ground, measured at 1-ft intervals, is given in the following table.
\begin{enumerate}[label=(\alph*)]
  \item Estimate the volume of the hole by using a Riemann sum with $m=n=3$ and the sample points to be the upper left corners of the subsquares of $R$.
  \item Approximate the average depth of the hole.\\ \begin{center} \begin{tabular}{lcccc} \hline  & \multicolumn{4}{c}{y} \\ x & 0 & 1 & 2 & 3 \\ 0 & 6 & 6.5 & 6.4 & 6 \\ 1 & 6.5 & 7 & 7.5 & 6.5 \\ 2 & 6.5 & 6.7 & 6.5 & 6 \\ 3 & 6 & 6.5 & 5 & 5.6 \\ \hline \end{tabular} \end{center}
\end{enumerate}
\item The level curves $f(x,y)=k$ of the function $f$ are given in the following graph, where $k$ is a constant.
\begin{enumerate}[label=(\alph*)]
  \item Apply the midpoint rule with $m=n=2$ to estimate the double integral $\underset{R}{\iint}f(x,y)\,dA$, where $R=[0.2,1]\times [0,0.8]$.
  \item Estimate the average value of the function $f$ on $R$.\\ \textit{[Figure: A series of curves marked $k=-1$, $-\frac{1}{2}$, $-\frac{1}{4}$, $-\frac{1}{8}$, $0$, $\frac{1}{8}$, $\frac{1}{4}$, $\frac{1}{2}$, and $1$. The line marked $k=0$ serves as an asymptote along the line $y=x$. The lines originate at $1$, $0.7$, $0.5$, $0.38$, and $0$ along the $y$-axis and at $0.38$, $0.5$, $0.7$, and $1$ along the $x$-axis, with the further-out lines curving less dramatically toward the asymptote.]}
\end{enumerate}
\item The level curves $f(x,y)=k$ of the function $f$ are given in the following graph, where $k$ is a constant.
\begin{enumerate}[label=(\alph*)]
  \item Apply the midpoint rule with $m=n=2$ to estimate the double integral $\underset{R}{\iint}f(x,y)\,dA$, where $R=[0.1,0.5]\times [0.1,0.5]$.
  \item Estimate the average value of the function $f$ on $R$.\\ \textit{[Figure: A series of quarter circles drawn in the first quadrant marked $k=\frac{1}{32}$, $\frac{1}{16}$, $\frac{1}{8}$, $\frac{1}{4}$, $\frac{1}{2}$, $\frac{3}{4}$, and $1$. The quarter circles have radii $0.17$, $0.25$, $0.35$, $0.5$, $0.71$, $0.87$, and $1$, respectively.]}
\end{enumerate}
\item The solid lying under the surface $z=\sqrt{4-y^{2}}$ and above the rectangular region $R=[0,2]\times [0,2]$ is illustrated in the following graph. Evaluate the double integral $\underset{R}{\iint}f(x,y)\,dA$, where $f(x,y)=\sqrt{4-y^{2}}$, by finding the volume of the corresponding solid. \\ \textit{[Figure: A quarter cylinder with center along the $x$-axis and with radius $2$. It has height $2$ as shown.]}
\item The solid lying under the plane $z=y+4$ and above the rectangular region $R=[0,2]\times [0,4]$ is illustrated in the following graph. Evaluate the double integral $\underset{R}{\iint}f(x,y)\,dA$, where $f(x,y)=y+4$, by finding the volume of the corresponding solid. \\ \textit{[Figure: In $xyz$-space, a shape is created with sides given by $y=0$, $x=0$, $y=4$, $x=2$, $z=0$, and the plane that runs from $z=4$ along the $y$-axis to $z=8$ along the plane formed by $y=4$.]}
\end{enumerate}

In the following exercises, calculate the integrals by interchanging the order of integration.

\begin{enumerate}[resume]
\item $\int_{-1}^{1}(\int_{-2}^{2}(2x+3y+5)\,dx)\,dy$
\item $\int_{0}^{2}(\int_{0}^{1}(x+2e^{y}-3)\,dx)\,dy$
\item $\int_{1}^{27}(\int_{1}^{2}(\sqrt[3]{x}+\sqrt[3]{y})\,dy)\,dx$
\item $\int_{1}^{16}(\int_{1}^{8}(\sqrt[4]{x}+2\sqrt[3]{y})\,dy)\,dx$
\item $\int_{\ln 2}^{\ln 3}(\int_{0}^{l}e^{x+y}\,dy)\,dx$
\item $\int_{0}^{2}(\int_{0}^{1}3^{x+y}\,dy)\,dx$
\item $\int_{1}^{6}(\int_{2}^{9}\frac{\sqrt{y}}{x^{2}}\,dy)\,dx$
\item $\int_{1}^{9}(\int_{4}^{2}\frac{\sqrt{x}}{y^{2}}\,dy)\,dx$
\end{enumerate}

In the following exercises, evaluate the iterated integrals by choosing the order of integration.

\begin{enumerate}[resume]
\item $\int_{0}^{\pi}\,\int_{0}^{\pi/2}\sin(2x)\cos(3y)\,dx\,dy$
\item $\int_{\pi/12}^{\pi/8}\,\int_{\pi/4}^{\pi/3}[\cot x+\tan(2y)]\,dx\,dy$
\item $\int_{1}^{e}\,\int_{1}^{e}[\frac{1}{x}\sin(\ln x)+\frac{1}{y}\cos(\ln y)]\,dx\,dy$
\item $\int_{1}^{e}\,\int_{1}^{e}\frac{\sin(\ln x)\cos(\ln y)}{xy}\,dx\,dy$
\item $\int_{1}^{2}\,\int_{1}^{2}(\frac{\ln y}{x}+\frac{x}{2y+1})\,dy\,dx$
\item $\int_{1}^{e}\,\int_{1}^{2}x^{2}\ln(x)\,dy\,dx$
\item $\int_{1}^{\sqrt{3}}\,\int_{1}^{2}y\arctan(\frac{1}{x})\,dy\,dx$
\item $\int_{0}^{1}\,\int_{0}^{1/2}(\arcsin x+\arcsin y)\,dy\,dx$
\item $\int_{0}^{1}\,\int_{1}^{2}xe^{x+4y}\,dy\,dx$
\item $\int_{1}^{2}\,\int_{0}^{1}xe^{x-y}\,dy\,dx$
\item $\int_{1}^{e}\,\int_{1}^{e}(\frac{\ln y}{\sqrt{y}}+\frac{\ln x}{\sqrt{x}})\,dy\,dx$
\item $\int_{1}^{e}\,\int_{1}^{e}(\frac{x\ln y}{\sqrt{y}}+\frac{y\ln x}{\sqrt{x}})\,dy\,dx$
\item $\int_{0}^{1}\,\int_{1}^{2}(\frac{x}{x^{2}+y^{2}})\,dy\,dx$
\item $\int_{0}^{1}\,\int_{1}^{2}\frac{y}{x+y^{2}}\,dy\,dx$
\end{enumerate}

In the following exercises, find the average value of the function over the given rectangles.

\begin{enumerate}[resume]
\item $f(x,y)=-x+2y$, $R=[0,1]\times [0,1]$
\item $f(x,y)=x^{4}+2y^{3}$, $R=[1,2]\times [2,3]$
\item $f(x,y)=\sinh x+\sinh y$, $R=[0,1]\times [0,2]$
\item $f(x,y)=\arctan(xy)$, $R=[0,1]\times [0,1]$
\item Let $f$ and $g$ be two continuous functions such that $0\le m_{1}\le f(x)\le M_{1}$ for any $x\in[a,b]$ and $0\le m_{2}\le g(y)\le M_{2}$ for any $y\in[c,d]$. Show that the following inequality is true: $m_{1}m_{2}(b-a)(d-c)\le \int_{a}^{b}\,\int_{c}^{d}f(x)g(y)\,dy\,dx\le M_{1}M_{2}(b-a)(d-c)$.
\end{enumerate}

In the following exercises, use property v. of double integrals and the answer from the preceding exercise to show that the following inequalities are true.

\begin{enumerate}[resume]
\item $\frac{1}{e^{2}}\le \underset{R}{\iint}e^{-x^{2}-y^{2}}\,dA\le1$, where $R=[0,1]\times [0,1]$
\item $\frac{\pi^{2}}{144}\le \underset{R}{\iint}\sin x\cos y\,dA\le \frac{\pi^{2}}{48}$, where $R=[\frac{\pi}{6},\frac{\pi}{3}]\times [\frac{\pi}{6},\frac{\pi}{3}]$
\item $0\le \underset{R}{\iint}e^{-y}\cos x\,dA\le \left( \frac{\pi}{2} \right)^{2}$, where $R=[0,\frac{\pi}{2}]\times [0,\frac{\pi}{2}]$
\item $0\le \underset{R}{\iint}(\ln x)(\ln y)\,dA\le{(e-1)}^{2}$, where $R=[1,e]\times [1,e]$
\item Let $f$ and $g$ be two continuous functions such that $0\le m_{1}\le f(x)\le M_{1}$ for any $x\in[a,b]$ and $0\le m_{2}\le g(y)\le M_{2}$ for any $y\in[c,d]$. Show that the following inequality is true: $(m_{1}+m_{2})(b-a)(d-c)\le \int_{a}^{b}\,\int_{c}^{d}[f(x)+g(y)]\,dy\,dx\le(M_{1}+M_{2})(b-a)(d-c)$.
\end{enumerate}

In the following exercises, use property v. of double integrals and the answer from the preceding exercise to show that the following inequalities are true.

\begin{enumerate}[resume]
\item $\frac{2}{e}\le \underset{R}{\iint}(e^{-x^{2}}+e^{-y^{2}})\,dA\le2$, where $R=[0,1]\times [0,1]$
\item $\frac{\pi^{2}}{36}\le \underset{R}{\iint}(\sin x+\cos y)\,dA\le \frac{\pi^{2}\sqrt{3}}{36}$, where $R=[\frac{\pi}{6},\frac{\pi}{3}]\times [\frac{\pi}{6},\frac{\pi}{3}]$
\item $\frac{\pi^{2}}{4}e^{\frac{-\pi}{2}}\le \underset{R}{\iint}(\cos x+e^{-y})\,dA\le \frac{\pi^{2}}{2}$, where $R=[0,\frac{\pi}{2}]\times [0,\frac{\pi}{2}]$
\item $0\le \underset{R}{\iint}(\ln x+\ln y)\,dA\le2\left( e-1 \right)^{2}$, where $R=[1,e]\times [1,e]$
\end{enumerate}

In the following exercises, the function $f$ is given in terms of double integrals.


(a) Determine the explicit form of the function $f$.


(b) Find the volume of the solid under the surface $z=f(x,y)$ and above the region $R$.


(c) Find the average value of the function $f$ on $R$.


(d) Use a computer algebra system (CAS) to plot $z=f(x,y)$ and $z=f_{\text{ave}}$ in the same system of coordinates.

\begin{enumerate}[resume]
\item \techrequired $f(x,y)=\int_{0}^{y}\,\int_{0}^{x}(xs+yt)\,ds\,dt$, where $(x,y)\in R=[0,1]\times [0,1]$
\item \techrequired $f(x,y)=\int_{0}^{x}\,\int_{0}^{y}[\cos(s)+\cos(t)]\,dt\,ds$, where $(x,y)\in R=[0,3]\times [0,3]$
\item Show that if $f$ and $g$ are continuous on $[a,b]$ and $[c,d]$, respectively, then $\int_{a}^{b}\,\int_{c}^{d}[f(x)+g(y)]\,dy\,dx=(d-c)\int_{a}^{b}f(x)\,dx+\int_{a}^{b}\,\int_{c}^{d}g(y)\,dy\,dx=(b-a)\int_{c}^{d}g(y)\,dy+\int_{c}^{d}\,\int_{a}^{b}f(x)\,dx\,dy$.
\item Show that $\int_{a}^{b}\,\int_{c}^{d}[yf(x)+xg(y)]\,dy\,dx=\frac{1}{2}(d^{2}-c^{2})\left(\int_{a}^{b}f(x)\,dx\right)+\frac{1}{2}(b^{2}-a^{2})\left(\int_{c}^{d}g(y)\,dy\right)$.
\item \techrequired Consider the function $f(x,y)=e^{-x^{2}-y^{2}}$, where $(x,y)\in R=[-1,1]\times [-1,1]$.
\begin{enumerate}[label=(\alph*)]
  \item Use the midpoint rule with $m=n=2,4,\ldots,\,10$ to estimate the double integral $I=\underset{R}{\iint}e^{-x^{2}-y^{2}}\,dA$. Round your answers to the nearest hundredths.
  \item For $m=n=2$, find the average value of $f$ over the region $R$. Round your answer to the nearest hundredths.
  \item Use a CAS to graph in the same coordinate system the solid whose volume is given by $\underset{R}{\iint}e^{-x^{2}-y^{2}}\,dA$ and the plane $z=f_{\text{ave}}$.
\end{enumerate}
\item \techrequired Consider the function $f(x,y)=\sin(x^{2})\cos(y^{2})$, where $(x,y)\in R=[-1,1]\times [-1,1]$.
\begin{enumerate}[label=(\alph*)]
  \item Use the midpoint rule with $m=n=2,4,\ldots,\,10$ to estimate the double integral $I=\underset{R}{\iint}\sin(x^{2})\cos(y^{2})\,dA$. Round your answers to the nearest hundredths.
  \item For $m=n=2$, find the average value of $f$ over the region $R$. Round your answer to the nearest hundredths.
  \item Use a CAS to graph in the same coordinate system the solid whose volume is given by $\underset{R}{\iint}\sin(x^{2})\cos(y^{2})\,dA$ and the plane $z=f_{\text{ave}}$.
\end{enumerate}
\end{enumerate}

In the following exercises, the functions $f_{n}$ are given, where $n\ge1$ is a natural number.


(a) Find the volume of the solids $S_{n}$ under the surfaces $z=f_{n}(x,y)$ and above the region $R$.


(b) Determine the limit of the volumes of the solids $S_{n}$ as $n$ increases without bound.

\begin{enumerate}[resume]
\item $f_{n}(x,y)=x^{n}+y^{n}+xy$, $(x,y)\in R=[0,1]\times [0,1]$
\item $f_{n}(x,y)=\frac{1}{x^{n}}+\frac{1}{y^{n}}$, $(x,y)\in R=[1,2]\times [1,2]$
\item Show that the average value of a function $f$ on a rectangular region $R=[a,b]\times [c,d]$ is $f_{\text{ave}}\approx \frac{1}{mn}\sum_{i=1}^{m}\sum_{j=1}^{n}f(x_{ij}^{*},y_{ij}^{*})$, where $(x_{ij}^{*},y_{ij}^{*})$ are the sample points of the partition of $R$, where $1\le i\le m$ and $1\le j\le n$.
\item Use the midpoint rule with $m=n$ to show that the average value of a function $f$ on a rectangular region $R=[a,b]\times [c,d]$ is approximated by \[f_{\text{ave}}\approx \frac{1}{n^{2}}\sum_{i=1}^{n}\sum_{j=1}^{n}f\left(\frac{x_{i-1}+x_{i}}{2},\frac{y_{j-1}+y_{j}}{2}\right).\]
\item An isotherm map is a chart connecting points having the same temperature at a given time for a given period of time. Use the preceding exercise and apply the midpoint rule with $m=n=2$ to find the average temperature over the region given in the following figure. \\ \textit{[Figure: A contour map showing surface temperature in degrees Fahrenheit. Given the map, the midpoint rule would give rectangles with values 71, 72, 40, and 43.]}
\end{enumerate}

\end{document}
